\printedpage{99}
% head: A. Grothendieck : Dérivateurs, Ch. XVIII, Édition des Universités de Montpellier II et Paris VII
% foot: Version 17 décembre 2025

a homomorphism on the $\varinjlim$ (\footnote{Careful, it is here that we tacitly use the fact that the functor $\lambda_* : \mathcal{E}_{j'} \longrightarrow \mathcal{E}_j$ commutes with $\varinjlim$ of type $I_\alpha$, which implies that the inclusions $\mathcal{E}_j \into \mathcal{E}$ commute with this type of $\varinjlim$ .}). Condition (i) ensures that $\lambda \longmapsto G(\lambda)$ transforms identities into identities, and condition (iii) that we have $G(\mu\lambda) = G(\mu)G(\lambda)$. Thus $G$ appears as a section of $\mathcal{E}$ over $J$. On the other hand, we have arrows
\[
X(j) \defeq X_0(j) \longrightarrow G(j) = \varinjlim_{\beta\in I_\alpha} X_\beta(j) \ ,
\]
which proves F (1). On the other hand, condition 3°) of the lemma proves the validity of F$'$ (1). This completes the proof of 4.10.3, and thereby of 4.10.1 and of 4.9.3.

\msp{page 263}

\textbf{Corollary 4.10.8} (of theorem 4.10.1). \textit{Let $\mathcal{N} \longrightarrow J$ be a cofibrant functor, such that the fibers are accessible categories (4.8), the cobase change functors $\lambda_*$ : $\mathcal{N}_{j'} \longrightarrow \mathcal{N}_j$ accessible, finally $J$ equivalent to a small category. Then the category $\mathcal{M} = \uGamma(\mathcal{N}/J)$ is accessible. More precisely, let $\pi_0$ be an infinite cardinal such that the categories $\mathcal{N}_j$ are $\pi_0$-accessible, and that the functors $\lambda_*$ are $\pi_0$-accessible, and let $\pi$ be an infinite cardinal such that we have:}
\begin{enumerate}[leftmargin=*]
\item[\cercle{a}] $\pi^{\pi_0} = \pi$.
\item[\cercle{b}] $\forall\ \lambda : j' \longrightarrow j$ \textit{in} $J$, $\lambda_* : \mathcal{E}_j \longrightarrow \mathcal{E}_{j'}$ \textit{maps} $(\mathcal{E}_j)_{\pi_0}$ \textit{into} $(\mathcal{E}_{j'})_{\pi_0}$.
\item[\cercle{c}] $\cardred J \leqslant \pi$.
\end{enumerate}
\textit{Then $\mathcal{M}$ is $\pi$-accessible, and more precisely, we have}
\[
\mathcal{M} \ \simeq \ \Ind_\pi(\mathcal{M}_\pi) \ ,
\]
\textit{where $\mathcal{M}_\pi$ is the full subcategory of $\mathcal{M}$ consisting of the sections $F$ of $\mathcal{M}$ over $J$ such that we have $F(j) \in (\mathcal{E}_j)_\pi \ \forall\ j \in \operatorname{Ob} J$,} \msp{page 264} \textit{$(\mathcal{E}_j)_\pi$ denoting the full subcategory of $\mathcal{E}_j$ consisting of the $\pi$-accessible objects of $\mathcal{E}_j$.}

Indeed, the hypothesis at the beginning implies the existence of $\pi_0$ and of $\pi$ satisfying the stated conditions. Conditions a) and b) on $\pi$, plus the fact that the functors $\lambda_*$ are $\pi_0$-accessible, imply that the strictly full subcategory $\mathcal{N}_\pi$ of $\mathcal{N}$ whose fibers are the $(\mathcal{N}_j)_\pi$, is a cofibered subcategory, and that moreover $\mathcal{N}$ is $J$-equivalent to $\Ind_\pi(\mathcal{N}_\pi/J)$. On the other hand, the fibers of $\mathcal{N}_\pi/J$ are stable under $\varinjlim$ large relative to $\pi_0$ and of cardinality $\leqslant \pi$, so the ``ordinal'' condition of 4.10.1 is satisfied (taking $\mathcal{D} = \mathcal{N}_\pi$ there). Since moreover $\operatorname{card} J \leqslant \pi$, we conclude by 4.10.1 that $\mathcal{M} \simeq \Ind_\pi(\mathcal{M}_\pi)$.

\msp{page 265}

\textbf{Corollary 4.10.9.} \textit{Let $\mathcal{N} \longrightarrow J$ be a ``cofibration in paratopoi'' (\footnote{ambiguous term, cf. the double explanation.}), i.e. a cofibrant functor such that the fibers are paratopos categories, and the cobase change functors are inverse image functors (resp. direct image functors) for morphisms of paratopoi, i.e. they commute with (small) $\varinjlim$ (resp. with small $\varprojlim$ ). Assume furthermore $J$ equivalent to a small category. Then $\uGamma(\mathcal{N}/J)$ is a paratopos category.}

Indeed, it is an accessible category by virtue of 4.10.8, on the other hand the hypothesis on the $\lambda_*$ implies at once that $\uGamma(\mathcal{N}/J)$ is stable under small $\varprojlim$ (resp. under small $\varprojlim$
\printedpage{100}
% head: A. Grothendieck : Dérivateurs, Ch. XVIII, Édition des Universités de Montpellier II et Paris VII
% foot: Version 17 décembre 2025
\texttt{[rather} $\varinjlim$ \texttt{?]}). Thus $\uGamma(\mathcal{N}/J)$ is a paratopos category by virtue of 4.8.9 b $\Longrightarrow$ a (resp. c $\Longrightarrow$ a).

\textbf{Corollary 4.10.10.} \textit{Let $\mathcal{N} \longrightarrow J$ be a category that is} fibered \textit{(not cofibered!), such that the fibers are paratopos categories and the functors} \msp{page 266} \textit{of base change are direct image functors for morphisms of paratopoi (i.e. accessible functors, commuting with small $\varprojlim$ ). Assume furthermore $J$ equivalent to a small category. Then $\uGamma(\mathcal{N}/J)$ is a paratopos category.}

Indeed, the hypothesis implies that for every arrow $\lambda : j' \longrightarrow j$ in $J$, the functor $\lambda_*$ : $\mathcal{N}_{j'} \longrightarrow \mathcal{N}_j$ admits a left adjoint $\lambda^* : \mathcal{N}_j \longrightarrow \mathcal{N}_{j'}$, whence it follows that $\mathcal{N} \longrightarrow J$ is also cofibrant, the cobase change functors commuting with small $\varinjlim$ . Thus we can apply 4.8.9, which gives the desired result.

\textbf{Question 4.10.11.} \textit{Let $\mathcal{N} \longrightarrow J$ be a functor that is} fibrant, \textit{such that the $\mathcal{N}_j$ are accessible categories, and the base change functors $\lambda^*$ accessible, $J$ being equivalent to a small category. I} \msp{page 267} \textit{conjecture that under these conditions, $\uGamma(\mathcal{N}/J) \defeq \Hom(J,\mathcal{N})$ is accessible, and that the same is true of $\Hom\,\underline{\operatorname{cart}}_J(J,\mathcal{N})$.}

This is what is proved in SGA 4 I 9.24 under slightly stronger hypotheses (with the help of the old notion too strong for my taste, of cardinal filtrations). I would be surprised if the (somewhat longish) proof of loc. cit. did not apply to the general case. But we would like to have a statement that covers this one and 4.10.8, in the case of a functor $\mathcal{N} \longrightarrow J$ that would be assumed neither fibrant nor cofibrant, but \ldots(???).

\section*{4.11\quad Application to the existence of limits in a category of the form $\Ind_{\pi_0}(C)$.}

\textbf{Proposition 4.11.1.} \textit{Let} \msp{page 268} \textit{$C$ be a category equivalent to a small category, $\pi$ an infinite cardinal, put $\mathcal{M} = \Ind_\pi(C)$. Let $\Phi$ be a set of categories of reduced cardinality $\leqslant \pi$. We propose to examine whether $\mathcal{M}$ is ``stable under $\varinjlim$ of type $\Phi$'' (i.e. stable under $\varinjlim$ of type $I$, whatever $I \in \Phi$ may be). It suffices for this that $C$ is stable under $\varinjlim$ of type $\Phi$, and that moreover there exists a limit ordinal $\alpha$ with $\operatorname{card}\alpha \leqslant \pi$ and such that $C$ is stable under $\varinjlim$ of type $I_\alpha$ (e.g. such that $I_\alpha \in \Phi$). Conversely, if $\mathcal{M}$ is stable under $\varinjlim$ of type $\Phi$, and if we assume moreover $C$ Karoubian (i.e. stable under direct factors), then $C$ is stable under $\varinjlim$ of type $\Phi$. In either case, the inclusion functor $C \into \mathcal{M}$ commutes with the $\varinjlim$ under consideration.}

(\footnote{Make explicit here the corollaries of types 4.13.20, 4.13.21 (p. 341).}).

\textsc{Proof}. a) Sufficient condition. We apply 4.9.3 to conclude that
\printedpage{101}
% head: A. Grothendieck : Dérivateurs, Ch. XVIII, Édition des Universités de Montpellier II et Paris VII
% foot: Version 17 décembre 2025
\msp{page 269}
for every category $J$ of reduced cardinality $\leqslant \pi$ (in particular for $J \in \Phi$), every functor $\varphi : J \longrightarrow \mathcal{M} = \Ind_\pi(C)$ admits an indicial representation $\Phi : J \times I \longrightarrow C$, with $I$ large relative to $\pi$. This, plus the stability of $C$ under $\varinjlim$ of type $J$, suffices to establish the existence of the $\varinjlim$ of $\varphi$.

The converse comes from the fact that if $C$ is Karoubian, then its essential image in $\mathcal{M}$ is none other than $\mathcal{M}_\pi$, which is stable in $\mathcal{M}$ under any $\varinjlim$ of cardinality $\leqslant \pi$ representable in $\mathcal{M}$. If $\mathcal{M}$ is stable under $\varinjlim$ of type $\Phi$, the same is true of $\mathcal{M}_\pi$, hence of $C$, q.e.d.

\textbf{Corollary 4.11.2.} \textit{For $\mathcal{M}$ to be stable under filtering $\varinjlim$ (resp. under small $\varinjlim$ ), it suffices that $C$ be so under filtering $\varinjlim$ (resp. under $\varinjlim$ )} of cardinality $\leqslant \pi$, \textit{and this is also necessary if $C$ is Karoubian.}

Suppose $C$ stable under filtering $\varinjlim$ \msp{page 270} (resp. under $\varinjlim$ of cardinality $\leqslant \pi$, then by 4.11.1 $\mathcal{M}$ is stable under this same type of direct limits. On the other hand, it is stable under $\varinjlim$ large relative to $\pi$. It thus suffices to prove the

\textbf{Lemma 4.11.3.} \textit{Let $\mathcal{M}$ be a category, $\pi$ a cardinal. Suppose that $\mathcal{M}$ is stable under filtering $\varinjlim$ of cardinality $\leqslant \pi$ (resp. under $\varinjlim$ of cardinality $\leqslant \pi$, resp. under sums of cardinality $\leqslant \pi$),} and \textit{under $\varinjlim$ large relative to $\pi$. Then $\mathcal{M}$ is stable under small filtering $\varinjlim$ (resp. under small $\varinjlim$ , resp. under small sums).} % [sic] des cardinal -> de cardinal, agreement (cf. "de cardinal" elsewhere in the statement)

This lemma is proved like 4.4.14, page 97.

There remains the converse in 4.11.2, which follows from the similar assertion in 4.11.1.

\textbf{Corollary 4.11.4.} \textit{When is $\mathcal{M} = \Ind_\pi C$ stable under sums? If there exists an ordinal $\alpha$ such that $\operatorname{card}\alpha \leqslant \pi$, and such that $C$ is stable under $\varinjlim$ of type $I_\alpha$, then it suffices that $C$ be stable under sums of cardinality $\leqslant \pi$. This condition is also necessary if $C$ is Karoubian.}

(\footnote{For other variants of 4.11.1, cf. 4.13.19 (p. 340) and corollaries.}).

Is proved like 4.11.2, using 4.11.3.

\msp{page 271}

\textbf{Proposition 4.11.5} (\footnote{?, proof not finished.}). \textit{Let $\Phi$ be a set of categories $I$ equivalent to finite categories, and ``rigid'', i.e. for every $i \in \operatorname{Ob} I$, $\Hom_I(i,i) = \{\mathrm{id}_i\}$. Let $C$ be a category equivalent to a small category, $\pi$ an infinite cardinal, $\mathcal{M} = \Ind_\pi C$. For $\mathcal{M}$ to be stable under $\varinjlim$ (resp. under $\varprojlim$ ) of type $\Phi$, it} suffices \textit{that $C$ be so. In the non-resp. case, if $C$ is Karoubian, this condition is also necessary.}

This last assertion is contained in the similar assertion in 4.11.1. As for the first assertion, this follows in the usual way from the

\textbf{Lemma 4.11.6} (\footnote{?, proof not spelled out is perhaps false!}). \textit{Let $J$ be a category equivalent to a finite category, $J$ rigid, $C$ a category equivalent to a small category, $\pi$ an infinite cardinal. Then every functor}
\printedpage{102}
% head: A. Grothendieck : Dérivateurs, Ch. XVIII, Édition des Universités de Montpellier II et Paris VII
% foot: Version 17 décembre 2025
\textit{$J \longrightarrow \mathcal{M} \defeq \Ind_\pi C$ admits an indicial representation $J \times I \longrightarrow C$, with $I$ large relative to $\pi$, in other words, the} \msp{page 272} \textit{canonical fully faithful functor}
\[
\Ind_\pi(\Hom(J,C)) \ \longrightarrow \ \Hom(J,\mathcal{M})
\]
\textit{deduced from the inclusion $C \into \mathcal{M}$, hence $\Hom(J,C) \into \Hom(J,\mathcal{M})$ (cf. 4.9), is an equivalence of categories.} (\textbf{N.B.} I do not know whether this is so without a rigidity hypothesis on $J$.)

\textsc{Proof}. This is proved like the similar result for $\Ind(C)$ instead of $\Ind_\pi C$, cf. SGA 4 I 8.8.5.

\msp{page 273}

% a vertical rule is printed in the left margin beside Corollaire 4.11.7
\textbf{Corollary 4.11.7} (\footnote{subject to the validity of 4.11.5.}). \textit{For $\mathcal{M} = \Ind_\pi C$ to be stable under $\varinjlim$ (resp. under $\varprojlim$ ) that are} finite, \textit{it} suffices \textit{that $C$ be so. And this is also necessary in the non-resp. case, when $C$ is assumed Karoubian.}

\section*{4.12\quad Accessible functors between accessible categories}

\textbf{Proposition 4.12.1} (\footnote{Would go back to 4.8.}). \textit{Let $\mathcal{M}$, $\mathcal{N}$ be two accessible categories (4.8), and $f : \mathcal{M} \longrightarrow \mathcal{N}$ an accessible functor. Then there exist infinite cardinals $\pi_0$, $\pi_1$ such that $\pi_1 \geqslant \pi_0$ and such that for every cardinal $\pi \geqslant \pi_1$, we have}
\[
f(\mathcal{M}_\pi) \ \subset \ \mathcal{N}_{\pi^{\pi_0}} \ , \tag*{(4.12.1)} % [sic] (4.12.1) -> (4.12.1.1), the formula is cited below as (4.12.1.1) and the next one is numbered (4.12.1.2)
\]
\textit{and in particular}
\[
f(\mathcal{M}_\pi) \ \subset \ \mathcal{N}_\pi \qquad \textit{si } \pi = \pi^{\pi_0}. \tag*{(4.12.1.2)}
\]
\textit{More precisely, we can take for $\pi_0$ a useful cardinal for $\mathcal{M}$} \msp{page 274} \textit{and such that $f$ is $\pi_0$-accessible (e.g. the smallest cardinal of this type). Then let $\pi_1$ be a cardinal such that we have $f(\mathcal{M}_{\pi_0}) \subset \mathcal{N}_{\pi_1}$, $\pi_1 \geqslant \cardred \mathcal{M}_{\pi_0}$, and finally such that $\mathcal{N}_{\pi_1}$ is} \texttt{[a]} \textit{strictly generating subcategory of $\mathcal{N}$. The pair $(\pi_0,\pi_1)$ satisfies the stated conditions.}

\textsc{Proof}. Indeed we will have
\[
\mathcal{N}_1 \ \xleftarrow{\ \sim\ } \Ind_{\pi_0}(\mathcal{M}_{\pi_0})
\]
and $f$, commuting with $\varinjlim$ large relative to $\pi_0$, will be known (up to unique isomorphism) by its restriction to $\mathcal{M}_{\pi_0}$, which factors through $\mathcal{N}_{\pi_1}$,
\[
f|\mathcal{M}_{\pi_0} : \mathcal{M}_{\pi_0} \ \longrightarrow \ \mathcal{N}_{\pi_1} \ .
\]
To prove (4.12.1.1), putting $\pi' = \pi^{\pi_0} \geqslant \pi$, and noting that $\pi'^{\pi_0} = \pi_0$ \texttt{[rather} $\pi'^{\pi_0} = \pi'$\texttt{]},
\[
f(\mathcal{M}_\pi) \ \subset \ f(\mathcal{M}_{\pi'}) \ \overset{?}{\subset} \ \mathcal{N}_{\pi'} \ ,
\]
\printedpage{103}
% head: A. Grothendieck : Dérivateurs, Ch. XVIII, Édition des Universités de Montpellier II et Paris VII
% foot: Version 17 décembre 2025
we see that it suffices to prove (4.12.1.2) (for $\pi'$). So we can assume $\pi = \pi^{\pi_0}$ and prove (4.12.1.2). But since $\pi_1 \geqslant \cardred \mathcal{M}_{\pi_0}$, we know that if $\pi$ is such that $\pi \geqslant \pi_1$, $\pi^{\pi_0} = \pi$, then $\mathcal{M}_\pi$ consists of the $\varinjlim\limits_{I} x_i$, the $x_i \in \mathcal{N}_{\pi_0}$, and $I$ large relative to $\pi_0$, and $\operatorname{card} I \leqslant \pi$.

\msp{page 275}

We then have ($f$ being $\pi_0$-accessible)
\[
f(x) \ = \ \varinjlim_{I} f(x_i) \ ,
\]
the $f(x_i)$ are in $\mathcal{N}_{\pi_1}$, hence in $\mathcal{N}_\pi$, now $\mathcal{N}_\pi$ stable under $\varinjlim$ of cardinality $\leqslant \pi$ (since $\mathcal{N}$ is generated by $\mathcal{N}_{\pi_1}$, consisting of objects that are $\pi_1$-accessible and a fortiori $\pi$-accessible, cf. SGA 4 I 9.9). Hence $f(x) \in \mathcal{N}_\pi$, q.e.d.

\textbf{Corollary 4.12.2.} \textit{Assume that $\pi_0$ is also useful for $\mathcal{N}$, and that $\pi_1 \geqslant \cardred \mathcal{N}_{\pi_0}$. Let $\pi$ be a cardinal $\geqslant \pi_1$, and such that $\pi^{\pi_0} = \pi$. Then $f$ is identified, up to equivalence of categories and isomorphism of functors, with the functor}
\[
\Ind(f_\pi) \ : \ \Ind_\pi(\mathcal{M}_\pi) \ \longrightarrow \ \Ind_\pi(\mathcal{N}_\pi) \ ,
\]
\textit{where $f_\pi : \mathcal{M}_\pi \longrightarrow \mathcal{N}_\pi$ is the functor induced by $f$.}

\textbf{Proposition 4.12.3} (\footnote{goes back to 4.7.}). \textit{Let $C$ be a category equivalent to a small category, $\mathcal{N}$ a $\mathfrak{U}$-category, $\pi_0$ an infinite cardinal.} \msp{page 276} \textit{We assume that $\mathcal{N}$ satisfies $L_{\pi_0}$ (stability under $\varinjlim$ large relative to $\pi_0$). Let $\mathcal{M} = \Ind_{\pi_0} C$, $\Hom_{!\,\pi_0}(\mathcal{M},\mathcal{N})$ the strictly full subcategory of $\Hom(\mathcal{M},\mathcal{N})$ consisting of the $f$ : $\mathcal{M} \longrightarrow \mathcal{N}$ that are $\pi_0$-accessible. Then the restriction functor}
\[
\Hom_{!\,\pi_0}(\mathcal{M},\mathcal{N}) \ \longrightarrow \ \Hom(C,\mathcal{N}) \tag*{(4.12.3.1)}
\]
\textit{is an equivalence of categories. We find a quasi-inverse functor as the composite}
\[
\Hom(C,\mathcal{N}) \ \longrightarrow \ \Hom(\underbrace{\Ind_{\pi_0}(C)}_{\mathcal{M}},\Ind_{\pi_0}(\mathcal{N})) \ \longrightarrow \ \Hom(\mathcal{M},\mathcal{N}) \ , \tag*{(4.12.3.2)}
\]
\textit{where the second arrow is the composition functor $F \longmapsto L \circ F$ with the direct limit functor $L : \Ind_{\pi_0}\mathcal{N} \longrightarrow \mathcal{N}$. The composite functor (4.11.3.2) has image contained in $\Hom_{!\,\pi_0}(\mathcal{M},\mathcal{N})$ and induces an equivalence between this last category and $\Hom(C,\mathcal{N})$.} % [sic] (4.11.3.2) -> (4.12.3.2), the composed functor is the display just labelled (4.12.3.2)

\msp{page 277}

\textit{In other words, we can say that the canonical inclusion functor}
\[
C \ \longrightarrow \ \Ind_{\pi_0} C \ = \ \mathcal{M}
\]
\textit{is} 2-universal \textit{for functors from $C$ into categories $\mathcal{N}$ satisfying $L_{\pi_0}$, the latter being regarded as $0$-objects of a $2$-category where the $1$-objects are the functors $\mathcal{N} \longrightarrow \mathcal{N}'$ that are $\pi_0$-accessible, i.e. commute with $\varinjlim$ large relative to $\pi_0$.}

This statement is the analogue, for arbitrary small $C$, of 4.4.16 (page 101), valid for $C$ stable under $\varinjlim$ of cardinality $\leqslant \pi_0$. There are of course other variants, in relation to commuting with other types of $\varinjlim$ ( filtering $\varinjlim$, sums, finite direct limits etc.) considered in the preceding section 4.11.

The proof of 4.12.3 is sorital.